\documentclass[11pt]{article}
\usepackage{explorer}
\newcommand{\ExplorerLanguage}{EN}
\newcommand{\ExplorerNumber}{03}
\newcommand{\ExplorerKind}{Lesson and method}
\begin{document}
\ExplorerTitle{High school · number theory}{Congruences and last digits}
\ExplorerSection{Before you begin}
Divisibility, powers and Euclidean division.\par
\ExplorerSection{Ideas and method}
\(a \equiv  b\) (mod m), with \(m \geq  2\), means that m divides \(a - b\). They have the same remainder upon division by m. For example, −\(1 \equiv  6\) (mod 7).\par
Congruences are preserved by addition and multiplication, hence by positive integer powers. Work modulo 10 to find a last digit.\par
A factor cannot always be cancelled: \(2 \times  1 \equiv  2 \times  4\) (mod 6), but 1 is not congruent to 4 modulo 6. A factor coprime to the modulus can be cancelled.\par
\ExplorerSection{Worked example}
The last digits of \(3^{1}, 3^{2}, 3^{3}, 3^{4}\) are 3,9,7,1. The cycle repeats because \(3^{4} \equiv  1\) (mod 10). Since \(2026 = 4 \times  506 + 2\), the last digit of \(3^{2026}\) is 9.\par
\ExplorerSection{Explore}
Build tables of squares modulo 3, 4 and 5. Identify impossible remainders.\par
\end{document}
